% !PRINTED 96
\weberpage{96}
%<<CONT
zeichen je ein volles Restsystem von $p$ (mit Ausschluss der 0) treten.

Wir wollen aber bei der Annahme $i_1=3$ stehen bleiben, und unter $\varepsilon_1, \varepsilon_2, \ldots \varepsilon_\mu$ irgend eine Combination der Zahlen $\pm 1$ verstehen. Dann können wir die Congruenz (3) auch so schreiben
\weq{4}{\varepsilon_1\alpha_1+\varepsilon_2\alpha_2+\cdots+\varepsilon_\mu\alpha_\mu\equiv 0\ (\mathrm{mod}\ 3).}

Die Zeichencombination $\varepsilon_1, \varepsilon_2, \ldots \varepsilon_\mu$, in der ein Zeichen willkürlich angenommen werden kann, da die gleichzeitige Aenderung aller Vorzeichen in (4) nichts ändert, bestimmt die einzelne Gruppe $\mathfrak{A}$, und offenbar sind diese Gruppen $\mathfrak{A}$ auch alle von einander verschieden. Denn ist z. B. in einer Gruppe $\varepsilon_1=\varepsilon_2=+1$, in einer anderen $\varepsilon_1=+1$, $\varepsilon_2=-1$, so enthält die erste die Zahlencombination:
\[
\alpha_1\equiv 1,\quad \alpha_2\equiv -1,\quad \alpha_3\equiv 0,\quad \ldots\ \alpha_\mu\equiv 0\ (\mathrm{mod}\ 3),
\]
die in der zweiten nicht vorkommt.

Um die Gruppen $\mathfrak{A}$ definitiv zu finden, legt man ein System primitiver Wurzeln $g_1, g_2, \ldots g_\mu$ von $q_1, q_2, \ldots q_\mu$ zu Grunde und bestimmt $a$ nach dem Modul $m$ aus den Congruenzen:
\weq{5}{\begin{aligned}
a &\equiv g_1^{\alpha_1}\ (\mathrm{mod}\ q_1)\\
&\equiv g_2^{\alpha_2}\ (\mathrm{mod}\ q_2)\\
&\cdots\cdots\cdots\cdots\cdots\\
&\equiv g_\mu^{\alpha_m}\ (\mathrm{mod}\ q_\mu). % [sic: Weber prints \alpha_m; the pattern requires \alpha_\mu.]
\end{aligned}}
Setzen wir, indem wir unter $\sigma$ einen der Reste 0, $\pm 1$ verstehen,
\weq{6}{\sigma \equiv \varepsilon_1\nu_1+\varepsilon_2\nu_2+\cdots+\varepsilon_\mu\nu_\mu\ (\mathrm{mod}\ 3),}
so besteht $\mathfrak{A}$ aus allen Zahlen $a$, deren Indices $\nu_1=\alpha_1, \nu_2=\alpha_2, \ldots \nu_\mu=\alpha_\mu$ der Bedingung $\sigma=0$ genügen. Die beiden anderen Werthe $\sigma=1$, $\sigma=-1$ bestimmen die beiden Nebengruppen
$\mathfrak{A}'$, $\mathfrak{A}''$ von $\mathfrak{A}$, so dass
\[
\mathfrak{N} = \mathfrak{A} + \mathfrak{A}' + \mathfrak{A}''
\]
ist. Ist $a'$ und $a''$ je eine Zahl aus $\mathfrak{A}'$ und $\mathfrak{A}''$, so ist $\mathfrak{A}'=\mathfrak{A}a'$, $\mathfrak{A}''=\mathfrak{A}a''$.

Wir bezeichnen nun mit $\eta$ die Kreistheilungsperiode, die zu $\mathfrak{A}$ gehört, und mit $\eta'$, $\eta''$ die conjugirten Perioden, also
\weq{7}{\eta = \sum r^a,\quad \eta' = \sum r^{a'},\quad \eta'' = \sum r^{a''}.}

Wenn dann $\varrho$ eine imaginäre dritte Einheitswurzel ist, so erhalten wir die Resolvente
\[
(\varrho, \eta) = \eta + \varrho\,\eta' + \varrho^2\,\eta'',
\]
%CONT>>
